\documentclass{article}

\usepackage[dblblindworkshop]{neurips_2026}

\usepackage[utf8]{inputenc}
\usepackage[T1]{fontenc}
\usepackage{hyperref}
\usepackage{url}
\usepackage{booktabs}
\usepackage{amsfonts}
\usepackage{amsmath}
\usepackage{nicefrac}
\usepackage{microtype}
\usepackage{xcolor}

\title{Positional or Relational? Syntactic Attention and Unmasking Order in
Diffusion Language Models}

\workshoptitle{Diffusion Language Models: Foundations, Efficiency, and Reasoning
(DiffuLM) at NeurIPS 2026}

\author{%
  Anonymous \\
}

\begin{document}

\maketitle

\begin{abstract}
Masked diffusion language models expose intermediate states in which most
positions are still \texttt{[MASK]}, making it possible to ask what a model has
committed to \emph{before} it commits to tokens. We propose to test two claims.
\textbf{H1:} dependency-tracking attention in masked DLMs separates into two
distinct head populations---positional heads implementing a fixed offset, and
relational heads that no offset can imitate---and the relational population
strengthens with scale. \textbf{H2:} that relational structure is not epiphenomenal:
it forms while both endpoints are masked, causally supports the eventual token,
and carries information about which positions the denoiser is ready to unmask,
making it a candidate signal for unmasking order. This proposal states both
hypotheses with explicit falsification criteria, the experiments that decide
them, and the infrastructure already in place.
\end{abstract}

\section{Background and motivation}

Diffusion LMs generate by iteratively denoising a fully masked sequence. Unlike
autoregressive models they expose intermediate states, so one can freeze the
model mid-generation and ask what structure exists while much of the sentence is
still \texttt{[MASK]}. Prior work by the authors\footnote{Non-archival workshop
paper; anonymized here per double-blind policy.} found attention heads in
DiffuGPT-S (12 layers $\times$ 12 heads) that predict the correct syntactic
receiver of a dependency, and found that the strongest such head is partially
correct \emph{before} either endpoint is revealed.

That work has one methodological flaw and one scope limit, and this proposal is
organized around fixing both.

\paragraph{The baseline was wrong.} A head that always attends a fixed number of
positions from itself solves any dependency whose endpoints sit at a
near-constant distance, while representing nothing syntactic. In English,
adjective$\to$noun and determiner$\to$noun endpoints are adjacent in most
instances, so a $+1$ head scores highly for free. The correct null is a
\emph{fixed-offset predictor} $\hat r = a + k$ with $k$ fit on the same split
used to select the head and reported held-out, so the baseline receives exactly
the tuning the head does. Under that null, only object$\to$verb and
subject$\to$verb survive in the small model (Table~\ref{tab:pilot}).

\paragraph{Why object$\to$verb is the informative case.} On full UD English-EWT,
the dependency-distance distributions are:

\begin{center}\small
\begin{tabular}{lccccc}
\toprule
relation & mode & 2nd & 3rd & 4th & best fixed offset \\
\midrule
object $\to$ verb        & $-2$: 35\% & $-1$: 34\% & $-3$: 17\% & $-4$: 6\% & 0.296 \\
subject $\to$ verb       & $+1$: 46\% & $+2$: 29\% & $+3$: 10\% & $+4$: 4\% & 0.461 \\
det.\ $\to$ noun         & $+1$: 57\% & $+2$: 28\% & $+3$: 8\%  & $+4$: 4\% & 0.594 \\
adj.\ $\to$ noun         & $+1$: 75\% & $+2$: 14\% & $+3$: 5\%  & $+4$: 2\% & 0.740 \\
\bottomrule
\end{tabular}
\end{center}

Object$\to$verb is the only relation with no dominant offset---its mass is split
almost evenly between $-1$ and $-2$ with a long tail---so a positional head is
capped near $0.30$ by construction. This is a property of the treebank, computed
with no model involved, and it is what makes the relation a usable probe.

\begin{table}[t]
\caption{Best head per relation against the fixed-offset null, recomputed from
prior runs. The small model clears the null on two relations; the 7B model
clears it on all six, and recruits a broad head population rather than one
lucky head. These are the numbers the proposed work must reproduce under a
corrected pipeline before anything is built on them.}
\label{tab:pilot}
\centering\small
\begin{tabular}{lrrrrrr}
\toprule
& \multicolumn{3}{c}{DiffuGPT-S (144 heads)} & \multicolumn{3}{c}{DiffuLLaMA-7B (1024 heads)} \\
\cmidrule(r){2-4}\cmidrule(l){5-7}
relation & null & head & $>$null & null & head & $>$null \\
\midrule
object $\to$ verb        & 0.314 & \textbf{0.755} & 29/144 & 0.277 & \textbf{0.877} & 151/1024 \\
subject $\to$ verb       & 0.429 & \textbf{0.521} & 6/144  & 0.433 & \textbf{0.660} & 46/1024 \\
obj.\ adj.\ $\to$ noun   & 0.750 & 0.827 & 1/144  & 0.750 & 0.942 & 13/1024 \\
subj.\ adj.\ $\to$ noun  & 0.682 & 0.682 & 2/144  & 0.682 & 0.886 & 19/1024 \\
obj.\ det.\ $\to$ noun   & 0.553 & 0.600 & 3/144  & 0.558 & 0.895 & 42/1024 \\
subj.\ det.\ $\to$ noun  & 0.600 & 0.589 & 3/144  & 0.600 & 0.856 & 30/1024 \\
\bottomrule
\end{tabular}
\end{table}

\section{Hypothesis 1: two head populations, separated by scale}

\paragraph{Statement.} Dependency-tracking attention in masked diffusion LMs is
carried by two distinguishable populations. \emph{Positional} heads implement a
fixed displacement: they solve every short-distance relation and fail on
object$\to$verb. \emph{Relational} heads solve object$\to$verb and score near
zero on the adjacent relations, a profile no fixed offset can produce. The
relational population is weak or absent at 124M parameters and clearly present
at 7B.

\paragraph{Preliminary evidence.} In the 7B pilot, two heads are near-exact
inverses of each other:

\begin{center}\small
\begin{tabular}{lccl}
\toprule
head & object $\to$ verb & the four adjacent relations & reading \\
\midrule
L18 H10 & 0.141 & 0.856 -- 0.942 & positional ($+1$ heuristic) \\
L3 H11  & 0.877 & 0.000 -- 0.023 & relational \\
\bottomrule
\end{tabular}
\end{center}

\paragraph{What would falsify it.} (i) The best object$\to$verb head at 7B also
solves the adjacent relations, i.e.\ no dissociation; (ii) accuracy as a
function of dependency distance decays identically for both head types, meaning
the ``relational'' head is a wider positional filter rather than a different
mechanism; (iii) the number of heads clearing the null is comparable across
relations, indicating the object$\to$verb result is a lucky draw from a large
search rather than a represented relation.

\paragraph{How it is decided.} The distance-stratified curve is the decisive
measurement, not the two hand-picked heads. A fixed-offset head's accuracy must
collapse once the dependency distance leaves its offset; a relational head's
must not. We bin held-out instances by $|$receiver $-$ attender$|$ and plot
accuracy per bin for the top head of each type, with the offset null on the same
axes. Two diagnostics accompany every relation and replace a multiple-comparisons
correction: how many of the searched heads clear the null at all, and the
Spearman rank correlation between selection-split and test-split accuracy across
all heads ($\rho \approx 0.97$ in the pilot).

\section{Hypothesis 2: the structure is used, and predicts unmasking order}

\paragraph{Statement.} Relational attention is not an artifact of the final,
fully-revealed state. It forms while both endpoints are still masked, it
causally supports the eventual token, and it carries information about which
positions the denoiser has resolved---making it a candidate control signal for
unmasking order, which is the free parameter every masked-DLM sampler must set.

\paragraph{Why this is the interesting claim for DiffuLM.} Masked diffusion
decoding chooses, at each step, which positions to reveal; the usual choices are
random, minimum entropy, or top-$k$ margin. All are \emph{token-confidence}
criteria. If the model resolves syntactic structure before lexical identity,
then a \emph{structural} criterion is available strictly earlier in the
trajectory than a confidence criterion, and prior measurement is consistent with
this: punctuation and function words are committed a median of seven diffusion
steps before they are unmasked, and earlier than content words (61\% / 54\% /
44\% committed early). H2 asks whether that ordering is real, exploitable, and
better than the token-confidence signals already in use.

\paragraph{What would falsify it.} (i) Masked-state accuracy does not exceed the
fixed-offset null computed \emph{on the same instances that enter the
masked-state average}---the pilot compared against a null fit on the whole
split, which is not a matched comparison; (ii) ablating the relational head over
the full held-out set does not reduce target-token support more than ablating
positional control heads matched on mean attention displacement, i.e.\ the effect
is generic loss of attention mass rather than loss of the relation; (iii)
relational-head state at step $t$ carries no information about which positions
the sampler reveals at $t{+}1$ beyond what token entropy already provides.

\paragraph{A caution we take seriously.} The pilot's masked-state accuracy is
$0.434$ at 124M but $0.108$ at 7B, i.e.\ \emph{below} the null. Two protocols
differ between those runs (five seeds with a $\ge$25-masked gate versus a single
run), so they are not comparable as they stand, and the metric itself may be
ill-posed: it scores the head against the gold parse of the \emph{original}
sentence, penalizing a model that denoises to a different but perfectly
grammatical sentence. Step~4 below fixes both before the number is interpreted.
If anticipation genuinely vanishes at scale, H2's first clause is false and we
will report that; it is a substantive result about diffusion LMs either way.

\section{Experimental plan}

Steps 1--3 establish H1, steps 4--6 test H2. Each produces an artifact that
stands alone if later steps fail.

\begin{enumerate}\itemsep2pt
\item \textbf{Corpus and null (done).} Full UD English-EWT, all splits pooled and
resampled at random under a fixed seed: 16{,}622 sentences $\to$ 9{,}016 usable
$\to$ disjoint 4000/1000/1000 splits, 21{,}434 relation instances. This is
4--5$\times$ the instances per relation of the pilot and, critically, the
selection split is now disjoint from the reporting split so head selection is
itself held out. The fixed-offset nulls move by at most a few points from the
pilot, which establishes that the earlier genre-ordered sampling did not distort
the baseline.

\item \textbf{Head search, both models.} Score all $L\times H$ heads on all six
relations at the final denoising frame, for DiffuGPT-S (144 heads) and
DiffuLLaMA-7B (1024 heads), on identical splits. $\approx$6{,}000 forward passes
per model. Deliverable: Table~\ref{tab:pilot} reproduced under the corrected
protocol, with Wilson intervals, and a per-head profile giving each head's
\emph{selectivity} (best relation minus second-best) and \emph{adjacency bias}
(mean over adjacent relations minus mean over distant ones). The two populations
of H1 should appear as separated clusters in that plane rather than as two
anecdotes.

\item \textbf{Distance-stratified curves.} For the top positional and top
relational head, accuracy against dependency distance with the offset null
overlaid. This is what converts H1 from ``two heads look different'' into a
mechanism claim.

\item \textbf{Masked-state, protocol-matched.} Re-run the timestep analysis
retaining \emph{per-instance} correctness, over five seeds, so the offset null
can be recomputed on exactly the instances entering each average. Trim the
schedule where the statistic is defined---the masked statistic dies once fewer
than 25 positions remain masked (step $\approx$38 of 64)---giving
$5 \times 21 \times 400 \approx 42$k forward passes rather than the 320k a naive
sweep would cost. Additionally re-score against a parse of the model's
\emph{own} completion, not the original sentence, to test whether the apparent
collapse at 7B is a metric artifact.

\item \textbf{Population-level causal ablation.} Zero the relational head across
the whole held-out set and measure the change in target-token logit, against
control heads matched on \emph{mean attention displacement} so the comparison
isolates the relation rather than the removal of attention mass. The prior
evidence was three hand-inspected sentences; this is the step that makes the
causal claim defensible.

\item \textbf{Unmasking-order probe (the payoff).} At each step, compare an
oracle reordering of the unmasking schedule driven by relational-head state
against the standard criteria (random, minimum entropy, top-$k$ margin) at a
fixed step budget, measured by likelihood of the reference and by agreement with
the model's own eventual output. We treat this first as a \emph{measurement}
(does the signal carry order-relevant information the confidence criteria lack?)
rather than as a decoder, because the measurement is cheap and answers the
scientific question; a competitive sampler is out of scope for this submission.
\end{enumerate}

\section{What already exists}

An installable package (\texttt{dlmrel}) implements steps 1--3 and the
scoring/aggregation for step 4, with the pipeline split so that everything not
needing a GPU---corpus construction, the null model, all analysis---runs and is
unit-tested first. It includes the fixed-offset null, the selection diagnostics,
Wilson and bootstrap intervals, deterministic split reconstruction across
processes, and a legacy-verification script that reproduces Table~\ref{tab:pilot}
from the prior runs' raw outputs, so a definitional drift cannot go unnoticed.
Steps 5 and 6 are not implemented.

Two failure modes are guarded at load time because both produce plausible
numbers silently: non-eager attention returns no attention weights while
accepting \texttt{output\_attentions=True}, and one of the public checkpoints is
serialized under a wrapper's parameter namespace, so loading it as a bare
backbone matches zero tensors and randomly initializes the network with only a
warning. The second cost us a full head search that came back at chance before
we caught it; the loader now refuses to proceed when a checkpoint populates less
than 90\% of the backbone.

\section{Risks}

\textbf{The 7B result does not reproduce.} The numbers in Table~\ref{tab:pilot}
come from a June run on genre-ordered sampling with a 64-token cap; they have
not been re-run. If the dissociation disappears under random full-treebank
sampling, H1 is false and the paper becomes a negative result about scale, which
we would report. \textbf{Compute.} Steps 2--5 are roughly two A100-hours; step 6
is the uncertain one and is deliberately scoped as a measurement.
\textbf{Metric validity.} Step 4's re-scoring against the model's own completion
may show the masked-state effect was never well-defined; that is the reason it
is a step rather than an assumption. \textbf{Scope fit.} DiffuLM's stated topics
are foundations, efficiency, and reasoning rather than interpretability, which is
why H2 is framed around unmasking order---the point at which this analysis
touches how diffusion LMs actually decode.

\end{document}
